\documentclass[11pt]{article}
\usepackage[a4paper,margin=18mm]{geometry}
\usepackage[no-math]{fontspec}
\setmainfont{Latin Modern Roman}
\usepackage{amsmath,amssymb,amsthm,xfrac,xspace,hyperref,graphicx,comment}
\excludecomment{DefTralics}
\providecommand{\C}{}
\providecommand{\bibliofont}{\small}
\newcommand{\ie}{i.e.,\xspace}
\newcommand{\eg}{e.g.,\xspace}
\newcommand{\etc}{etc.\xspace}
\newcommand{\cf}{cf.\xspace}
\newcommand{\resp}{resp.\xspace}
\newcommand{\smaller}{\small}
\newcommand{\Subsection}{\subsection}
\newcommand{\Subsubsection}{\subsubsection}
\newcommand{\skpt}{\smallskip}
\emergencystretch=3em
\providecommand{\dpl}{\displaystyle}
\input{author-preamble.tex}

\begin{document}
\begin{center}{\Large Stable item 1931}\end{center}
\paragraph{Source and attribution.}
Karine Beauchard, Armand Koenig, Kévin Le Balc’h, \emph{Null-controllability of linear parabolic-transport systems}, Journal de l’École polytechnique — Mathématiques 7 (2020), publisher version, DOI 10.5802/jep.127. \href{https://jep.centre-mersenne.org/articles/10.5802/jep.127/}{Publisher article}. Authors' text: \href{https://creativecommons.org/licenses/by/4.0/}{CC BY 4.0}.
\paragraph{Editorial problem boundary.}
Prove only Theorem 2(i) below for a nonempty proper open interval \(\omega\) of the circle, with all printed constant-matrix assumptions and full distributed control \(M=I_d\). For every finite \(0<T<T^*\), where \(T^*=(2\pi-|\omega|)/\mu_*\) when \(\mu_*>0\) and \(T^*=+\infty\) otherwise, the system is not null controllable. The perturbation/spectral decomposition, adjoint observability equivalence, polynomial-filtered compact transport packets, coupled-packet comparison and complete positive-final-norm contradiction are retained. The borderline time and all positive-control assertions are outside the selected conclusion.
\paragraph{Difficulty (editorial): H2.}
Holomorphic spectral projections isolate a genuinely hyperbolic branch of the coupled dissipative Fourier symbol. Support-preserving high-frequency polynomial cutoffs create exact adjoint solutions within order inverse-frequency of unobserved transport packets, disproving observability despite instantaneous parabolic coupling.
\paragraph{Editorial transcription note.}
Original author language, mathematical content and ordering are unchanged. Publisher front matter is replaced by this independent layout wrapper. Original native statement, proof and macros remain editable; bibliography is recovered from the matching cached publisher HTML. No assistant-generated proof or mathematical repair is supplied.
\clearpage
\input{author-body.tex}
\begin{thebibliography}{99}
\bibitem{alabau-boussouira_2017} F. Alabau-Boussouira, J.-M. Coron \& G. Olive - “Internal controllability of first order quasi-linear hyperbolic systems with a reduced number of controls” , SIAM J. Control Optim. 55 (2017) no. 1, p. 300-323
\bibitem{albano_2000} P. Albano \& D. Tataru - “Carleman estimates and boundary observability for a coupled parabolic-hyperbolic system” , Electron. J. Differential Equations (2000), article ID 22, 15 pages
\bibitem{ammarkhodja_2009a} F. Ammar Khodja, A. Benabdallah, C. Dupaix \& M. González-Burgos - “A generalization of the Kalman rank condition for time-dependent coupled linear parabolic systems” , Differential Equations Appl. 1 (2009) no. 3, p. 427-457
\bibitem{ammarkhodja_2009} F. Ammar Khodja, A. Benabdallah, C. Dupaix \& M. González-Burgos - “A Kalman rank condition for the localized distributed controllability of a class of linear parbolic systems” , J. Evol. Equ. 9 (2009) no. 2, p. 267-291
\bibitem{duprez_2016a} F. Ammar Khodja, F. Chouly \& M. Duprez - “Partial null controllability of parabolic linear systems” , Math. Control Relat. Fields 6 (2016) no. 2, p. 185-216
\bibitem{bardos_1992} C. Bardos, G. Lebeau \& J. Rauch - “Sharp sufficient conditions for the observation, control, and stabilization of waves from the boundary” , SIAM J. Control Optim. 30 (1992) no. 5, p. 1024-1065
\bibitem{beauchard_2011} K. Beauchard \& E. Zuazua - “Large time asymptotics for partially dissipative hyperbolic systems” , Arch. Rational Mech. Anal. 199 (2011) no. 1, p. 177-227
\bibitem{brezis_2011} H. Brezis - Functional analysis, Sobolev spaces and partial differential equations , Universitext , Springer, New York, 2011
\bibitem{chaves-silva_2014} F. W. Chaves-Silva, L. Rosier \& E. Zuazua - “Null controllability of a system of viscoelasticity with a moving control” , J. Math. Pures Appl. (9) 101 (2014) no. 2, p. 198-222
\bibitem{chowdhury_2015} S. Chowdhury \& D. Mitra - “Null controllability of the linearized compressible Navier-Stokes equations using moment method” , J. Evol. Equ. 15 (2015) no. 2, p. 331-360
\bibitem{chowdhury_2014} S. Chowdhury, D. Mitra, M. Ramaswamy \& M. Renardy - “Null controllability of the linearized compressible Navier-Stokes system in one dimension” , J. Differential Equations 257 (2014) no. 10, p. 3813-3849
\bibitem{coron_2007} J.-M. Coron - Control and nonlinearity , Math. Surveys and Monographs , vol. 136 , American Mathematical Society, Providence, RI, 2007
\bibitem{duprez_2018b} M. Duprez \& G. Olive - “Compact perturbations of controlled systems” , Math. Control Relat. Fields 8 (2018) no. 2, p. 397-410
\bibitem{ervedoza_2012} S. Ervedoza, O. Glass, S. Guerrero \& J.-P. Puel - “Local exact controllability for the one-dimensional compressible Navier-Stokes equation” , Arch. Rational Mech. Anal. 206 (2012) no. 1, p. 189-238
\bibitem{guerrero_2013} S. Guerrero \& O. Y. Imanuvilov - “Remarks on non controllability of the heat equation with memory” , ESAIM Control Optim. Calc. Var. 19 (2013) no. 1, p. 288-300
\bibitem{guzman_2019} P. Guzmán \& L. Rosier - “Null controllability of the structurally damped wave equation on the two dimensional torus” , 2019
\bibitem{ivanov_2009} S. Ivanov \& L. Pandolfi - “Heat equation with memory: lack of controllability to rest” , J. Math. Anal. Appl. 355 (2009) no. 1, p. 1-11
\bibitem{kato_1995} T. Kato - Perturbation theory for linear operators , Classics in Math. , vol. 132 , Springer, Berlin Heidelberg, 1995
\bibitem{koenig_2017} A. Koenig - “Non-null-controllability of the Grushin operator in 2D” , Comptes Rendus Mathématique 355 (2017) no. 12, p. 1215-1235
\bibitem{lerousseau_2012} J. Le Rousseau \& G. Lebeau - “On Carleman estimates for elliptic and parabolic operators. Applications to unique continuation and control of parabolic equations” , ESAIM Control Optim. Calc. Var. 18 (2012) no. 3, p. 712-747
\bibitem{lebeau_1995} G. Lebeau \& L. Robbiano - “Contrôle exact de l’équation de la chaleur” , Comm. Partial Differential Equations 20 (1995-01) no. 1-2, p. 335-356
\bibitem{lebeau_1998} G. Lebeau \& E. Zuazua - “Null-controllability of a system of linear thermoelasticity” , Arch. Rational Mech. Anal. 141 (1998) no. 4, p. 297-329
\bibitem{martin_2013} P. Martin, L. Rosier \& P. Rouchon - “Null controllability of the structurally damped wave equation with moving control” , SIAM J. Control Optim. 51 (2013) no. 1, p. 660-684
\bibitem{peetre_1961} J. Peetre - “Another approach to elliptic boundary problems” , Comm. Pure Appl. Math. 14 (1961) no. 4, p. 711-731
\bibitem{queffelec_2013} H. Queffélec \& C. Zuily - Analyse pour l’agrégation , Dunod, Paris, 2013
\bibitem{rosier_2007} L. Rosier \& P. Rouchon - “On the controllability of a wave equation with structural damping” , Int. J. Tomogr. Stat. 5 (2007) no. W07, p. 79-84
\bibitem{tucsnak_2009} M. Tucsnak \& G. Weiss - Observation and control for operator semigroups , Birkhäuser Advanced Texts , Birkhäuser Verlag, Basel, 2009
\end{thebibliography}
\end{document}
